\chapter{Phase--parameter correspondence}

\section{External-ray data}

The normalized Böttcher coordinate gives a potential--angle
parametrization of the basin of infinity.  In the formalization we separate
the analytic parametrization from its boundary behavior.

\begin{definition}[External-ray system]
\label{def:external-ray-system}
For a quadratic map $f_c$, an external-ray system consists of maps
\[
 \gamma_{c,\theta}:[0,\infty)\to\C,\qquad
 \ell_c(\theta)\in\C,\qquad\theta\in\R/\Z,
\]
such that
\[
 \lim_{t\to+\infty}\gamma_{c,\theta}(t)=\infty,\qquad
 \lim_{t\to0^+}\gamma_{c,\theta}(t)=\ell_c(\theta)
\]
whenever the ray lands.  The dynamical relation is
\[
 f_c(\gamma_{c,\theta}(t))
 =\gamma_{c,2\theta}(2t).
\]
The landing map $\ell_c$ is allowed to be partial in the analytic
construction; a finite puzzle graph uses only the angles for which landing
has been proved.
\end{definition}

\begin{remark}
The Lean structure \lean{ExternalRaySystem} records a total landing map as
an explicit contract.  This is intentionally stronger than merely defining
an external ray: the boundary limit is the information needed to form a
closed graph and to transport combinatorics to parameter space.
\end{remark}

\section{Dynamical puzzle graphs}

\begin{definition}[Finite dynamical graph]
\label{def:finite-dynamical-graph}
Choose finitely many angles whose rays land and add the equipotential arcs
of a fixed positive potential.  Pulling this finite graph back under
$f_c^{\circ n}$ produces a graph $\Gamma_n(c)$.  The critical puzzle piece
is the component
\[
 P_n(c):=\operatorname{Comp}_0(\C\setminus\Gamma_n(c)).
\]
\end{definition}

\begin{proposition}[Puzzle properties]
\label{prop:puzzle-properties}
The required finite-level data are:
\begin{enumerate}
\item $\Gamma_n(c)$ is closed and $P_n(c)$ is open;
\item $0\in P_n(c)$ and $P_n(c)$ is connected;
\item $P_{n+1}(c)\subseteq P_n(c)$;
\item the annulus $P_n(c)\setminus\overline{P_{n+1}(c)}$ carries a
      conformal modulus.
\end{enumerate}
\end{proposition}

The first three properties are topological.  The fourth is analytic and
requires the graph to avoid critical degeneracies.  In Lean they are fields
of \lean{DynamicalPuzzleSystem}; the structure does not assert that such a
system exists for every parameter.

\section{Parameter graphs}

\begin{definition}[Phase--parameter realization]
\label{def:phase-parameter-realization}
A phase--parameter realization at $c\in\M$ consists of parameter graphs
$\mathcal G_n$ and parameter pieces $\mathcal P_n(c)$ such that:
\[
\begin{array}{ll}
\mathcal G_n\text{ is closed},&
\mathcal P_n(c)\text{ is open},\\
\mathcal P_{n+1}(c)\subseteq\mathcal P_n(c),&
c\in\mathcal P_n(c),\\
\mathcal P_n(c)\cap\M\text{ is connected},&
\bigcap_n\mathcal P_n(c)=\{c\}.
\end{array}
\]
The phase--parameter compatibility additionally identifies the finite
combinatorial data of $\mathcal P_n(c)$ with those of $P_n(c)$.
\end{definition}

\begin{theorem}[Topological transport]
\label{thm:topological-transport}
If the parameter pieces in Definition~\ref{def:phase-parameter-realization}
form a neighbourhood basis at $c$ in the parameter plane, then $\M$ is
locally connected at $c$.
\end{theorem}

\begin{proof}
The connected sets $\mathcal P_n(c)\cap\M$ are neighbourhoods of $c$ in
the subspace $\M$ and refine every neighbourhood of $c$.  They therefore
form a connected neighbourhood basis.  The Lean theorem
\lean{localConnectivity_of_phaseParameterData} proves exactly this
topological implication from \lean{PhaseParameterPuzzleData}.
\end{proof}

\begin{conditionalstep}
The difficult theorem is not Theorem~\ref{thm:topological-transport}, but
the existence of the data in Definition~\ref{def:phase-parameter-realization}.
It combines external-ray landing, graph pullbacks, holomorphic motion,
properness of the phase--parameter map, and the identification of the
parameter image.  In the review this is the explicit proposition
\lean{YoccozPhaseParameterResult}; it is not hidden in a global MLC axiom.
\end{conditionalstep}

\section{The full Green-sublevel warning}

For $c\in\M$ one can form the larger sets
\[
 \widetilde P_n(c):=
 \{c'\in\C:G_c(c'-c)<2^{-n}\}.
\]
They are useful analytic neighbourhoods, but they are not automatically
the graph-cut parapuzzle pieces of Definition~\ref{def:phase-parameter-realization}.
In fact, the full-sublevel intersection can retain a translate of the
filled Julia set.  Thus the implication
\[
 \widetilde P_n(c)\text{ connected}
 \quad\Longrightarrow\quad
 \mathcal P_n(c)\text{ is a valid parapuzzle}
\]
is invalid without a separate carving or quotient theorem.

\begin{openestimate}[Phase--parameter existence]
\label{est:phase-parameter}
Construct the graph-cut phase--parameter realization for the precise
combinatorial classes used by the renormalization proofs, and prove that
its parameter pieces are connected, nested, and shrinking.
\end{openestimate}
